具有相当复杂程度的数学对象，常常伴有更高阶的\trHomotopy{}（homotopy）结构：这些
对象之间的映射可以由见证图之弱交换性的\trHomotopy{}相连；这些\trHomotopy{}又可以由表达诸
见证之间相干性条件（coherence conditions）的更高阶\trHomotopy{}相连；而这些更高阶
的\trHomotopy{}还可以由再高一阶的\trHomotopy{}相连，以至无穷。这类数学对象的自然栖居之所
并非通常的 1-\trCategory{}，而是\textbf{\trInfinityCategory}（∞-category），或更
精确地说，$(\infty,1)$-\trCategory{}（$(\infty,1)$-category）；指标「1」意指：
最低维以上——即刚才讨论的那些\trHomotopy——的态射（morphisms）都是弱可逆的。

这里，定义\trInfinityCategory{}之高阶态射的\trHomotopy{}，应被当作\emph{数据}（data）来看待，
而不只是 1 维态射所承载的\trEquivalence{}关系（equivalence relation）的单纯见证。
这一视角的转变，可以由拓扑空间（topological space）的两个代数不变量之间
的关系来例示：基本\trGroupoid{}（fundamental groupoid）是一个通常的 1-\trCategory{}；基本%
\trInfinityGroupoid{}（fundamental ∞-groupoid）则是一个所有态射皆弱可逆的\trInfinityCategory{}。两种
情形下，对象都是外围拓扑空间的点；但在前者中，1-态射是道路的\trHomotopy{}类，而在
后者中，1-态射是道路本身，2-态射则是保持端点的显式\trHomotopy{}。为了容纳这类例子，
\trInfinityCategory{}中的一切\trCategory{}性构造都是弱的。即便在 1-态射这一基层上，合成
（composition）也未必唯一确定：它由一个 2-态射所见证，并且在 3-态射之下
结合，而这个 3-态射的边界数据又涉及指定的 2-态射见证。因此，取值于\trInfinityCategory
中的图（diagram）不能说\emph{不折不扣地交换}，而应诠释为\emph{\trHomotopy{}相干}
的，并带有显式给定的高阶数据。

定义\trInfinityCategory{}的一项根本性挑战，在于为这种弱合成律（weak composition law）
的观念给出精确的数学含义——不仅对 1-态射，也要对更高维度的态射。的确，
在某种意义上，我们传统的以集合为基础的数学基础并不真正适合推理\trInfinityCategory{}：
集合在\trInfinityCategory{}式的数据中并不显要，尤其当我们以道德上正确的方式对待\trInfinityCategory{}，
把它们视为仅在\textbf{\trEquivalence}（equivalence）意义下良好定义的对象
时，更是如此。在\trEquivalence{}意义下考量时，\trInfinityCategory{}与通常的\textbf{\trCategory}
（category）一样，并没有一个良好定义的对象集。此外，\trInfinityCategory{}中一对固定对象
之间的态射装配成一个\trInfinityGroupoid{}，后者描述了一个良好定义的\trHomotopy{}型（homotopy
type），尽管并非良好定义的空间。\footnote{Grothendieck 的\trHomotopy{}假设主张：
\trInfinityGroupoid{}在\trEquivalence{}意义下对应于\trHomotopy{}型。}

精确性是借助 $(\infty,1)$-\trCategory{}的种种\emph{模型}（model）达成的：它们是
Bourbaki 式的数学结构，用以表示带有弱合成律的无穷维\trCategory{}，其中维数 1 以上
的态射皆弱可逆。按出场顺序，这些模型包括：\trSimplicialCategory{}（simplicial
categories）、\textbf{\trQuasiCategory}（quasi-category，旧称「弱\trKanComplex{}」）、
相对\trCategory{}（relative categories）、\trSegalCategories{}（Segal categories）、
\trCompleteSegalSpaces{}（complete Segal spaces），以及\trOneComplicialSets{}（1-complicial
sets，旧称「饱和 1-平凡弱复状集」），每一种都伴有一大批自然出现的实例。
$(\infty,1)$-\trCategory{}模型的这种繁衍，引出了它们该如何比较的问题。二十一世纪的
最初几十年里，Julia Bergner、André Joyal 与 Myles Tierney、Dominic Verity、Jacob Lurie，以及 Clark Barwick 与 Daniel Kan 架起了种种桥梁，
证明上列每一种模型都「具有相同的\trHomotopy{}论」——其含义是：它们分别给出 Quillen
\trEquivalence{}的\textbf{\trModelCategory}（model category）中的\trFibrant{}对象（fibrant
objects）。\footnote{Bergner 近年的一部专著综述了上述模型中除最后一个之外
的全部模型及其相互关系 [15]；若要更走马观花式的巡览，可见 [3]。}

在 $(\infty,1)$-\trCategory{}诸模型的发展与模型间比较构造的推进并行之际，一项
大获成功的计划也在铺展：Joyal 开创了它，Lurie 与众多后来者则将其发扬光大
——把基础范畴论（category theory）从通常的 1-\trCategory\trExtension{}到以\trQuasiCategory{}为模型的
$(\infty,1)$-\trCategory{}，使得新的\trQuasiCategory{}式概念沿标准嵌入
$\mathcal{C}at\hookrightarrow\mathcal{Q}\mathcal{C}at$ 限制时给出经典的
1-\trCategory{}概念。于是自然要问：这项工作能否\trExtension{}到 $(\infty,1)$-\trCategory{}的其他模型？
基础的\trInfinityCategory{}式概念在\trChangeOfModel{}之下又在多大程度上保持不变？例如，流形
（manifold）的 $(\infty,1)$-\trCategory{}最自然地被构造为\trCompleteSegalSpaces{}，因此
Kazhdan–Varshavsky [65]、Boavida de Brito [34] 与 Rasekh [95, 96, 98]
近来致力于以\trCompleteSegalSpaces{}重新发展\trQuasiCategory{}的范畴论中的一些内容，以便直接
取得那些此前仅在\trQuasiCategory{}模型中引入的构造与定义。

出于实际、审美与道德的理由，实践者的终极愿望是「模型无关地」（model
independently）工作：用 $(\infty,1)$-\trCategory{}的任一模型证明的定理对它们全体
成立，而任何特定模型固有的技术细节从不进入讨论。既然 $(\infty,1)$-\trCategory{}的
所有模型都「具有相同的\trHomotopy{}论」，普遍共识是模型的选择不应有太大影响；但
证明这类结果的一个障碍是：已为\trQuasiCategory{}建立的种种\trCategory{}定义之精确版本，在很大
程度上尚未对其他模型给出。在不同模型之间确有可比定义的情形，所论\trCategory{}概念
的模型不变性（model invariance）通常也能给出一个 ad hoc 的启发性证明，
细节则留给精通各模型组合学的专家去补足；但若有一种系统的方法，能借助多少
更贴近地面的论证去比较不同模型中 $(\infty,1)$-\trCategory{}的范畴论，那会令人踏实
得多。
